\documentclass[11pt,a4paper]{ctexart}
\usepackage[margin=2.4cm]{geometry}
\usepackage{amsmath,amssymb,bm}
\usepackage{amsthm}
\usepackage{enumitem}
\usepackage{fancyhdr}
\usepackage[hidelinks]{hyperref}
\pagestyle{fancy}
\fancyhf{}
\fancyhead[L]{Team 2608797}
\fancyhead[R]{\thepage}
\renewcommand{\headrulewidth}{0.4pt}
\setlength{\parskip}{0.35em}
\setlength{\parindent}{2em}
\linespread{1.15}
\newtheorem{theorem}{定理}
\newtheorem{definition}{定义}
\newtheorem{proposition}{命题}
\newtheorem{lemma}{引理}
\renewcommand{\proofname}{证明}
\newcommand{\degC}{\ensuremath{^{\circ}\mathrm{C}}}
\setlist[itemize]{topsep=0.25em,itemsep=0.15em,parsep=0.1em}
\setlist[enumerate]{topsep=0.25em,itemsep=0.15em,parsep=0.1em}
\title{多尺度电热耦合的智能手机电池续航预测模型}
\author{}
\date{}
\begin{document}
\maketitle

\begin{center}
\small 选题：A \quad 2026 MCM/ICM 摘要页 \quad 队伍控制号：2608797
\end{center}
\vspace{0.5em}

\section*{摘要}
\addcontentsline{toc}{section}{摘要}

智能手机电池的不可预测性始终是用户与工程师共同面临的重大挑战。本文构建了一个完整的连续时间状态空间模型（Continuous-Time State-Space Model），用以刻画荷电状态（State of Charge, SOC）的动态演化，并在复杂、真实的使用场景下准确预测耗尽时间（Time-to-Empty, TTE）。

我们的建模框架融合了电化学动力学、热反馈与长期老化动力学。其核心采用二阶 Thevenin 等效电路（RC 网络）以捕捉非线性电压恢复与极化效应。该模型通过将电池内阻与基于 Arrhenius 的热模块、基于 SEI 生长的老化模块相耦合，解决了“电压—电流代数环”问题。为弥合离散用户行为与连续动力学之间的鸿沟，我们采用屏幕亮度的 Gamma 非线性模型与网络信号强度的双曲补偿模型来模拟功耗。此外，随机使用模式通过用于亮度波动的 Ornstein--Uhlenbeck（O-U）过程与用于网络“尾部能量”事件的 Bernoulli 过程引入。

跨多种基准的数值仿真验证了模型的高保真度：在正常浏览模式下，与实测数据相比预测误差仅为 0.5\%。关键发现表明，在极端寒冷（$-10\,\degC$）条件下，设备会在仍保有 38.4\% SOC 时经历“电压崩溃”关机，我们定量地将这一现象归因于内阻的分层 Arrhenius 响应。采用 Morris 筛选法的敏感性分析识别出欧姆活化能在寒冷环境中是“杀手参数”，而在室温下行为随机性主导 TTE 方差（92\%）。我们进一步证明，由于基线电阻漂移与热应力的协同作用，老化电池在寒冷场景下的续航比新电池缩短 35\%。

基于上述洞见，我们提出一种物理状态感知的优化策略。对用户，我们建议将亮度上限设为 70\%--80\% 以避开“电池自杀”区，并在弱信号区域启用飞行模式以缓解双曲型功率激增。对操作系统，我们建议采用阻抗感知的 VS-DVFS 策略与主动热拦截以防止泄漏诱发的失控。最后，我们通过将该模型成功适配至无人机（Unmanned Aerial Vehicles, UAVs）与可穿戴设备，证明了其可扩展性，确立了其作为便携式能源管理通用框架的稳健性。

\noindent\textbf{关键词：}连续时间状态空间表示；二阶 Thevenin 电路；热—老化耦合；Ornstein--Uhlenbeck 过程；耗尽时间预测

\newpage
\section{引言}

\subsection{问题背景}

在移动互联网时代，智能手机已成为人们生活中的必需品。然而，续航始终是用户的核心痛点。电池性能受到高度复杂且动态的因素影响。作为动力核心，锂离子电池的能量释放并非简单的线性消耗过程，而是由用户行为、设备运行状态与外部环境条件多变量耦合所支配的电化学反应。预测“耗尽时间”（Time-to-Empty, TTE）的挑战在于电池的非线性特性——例如电压跌落、极化效应~[8] 与温度敏感性——这些特性往往导致实际可用容量与标称电量之间出现显著偏差。精确预测电池在不同使用模式下的剩余时间并识别关键影响因素，对于缓解“电池焦虑”、指导系统优化与硬件设计具有重要意义。

\subsection{问题重述}

本研究的核心目标是建立一个连续时间数学模型，以准确刻画并预测智能手机在真实使用场景下的电池动态行为。具体目标如下：

\begin{itemize}
\item \textbf{荷电状态（SOC）动态演化}：我们旨在建立连续函数 $\mathrm{SOC}(t)$，其中时间 $t$ 为自变量。这将通过一组常微分方程（ODEs）实现，用以描述不同负载条件下的能量耗尽速率~[1]。
\item \textbf{多源负载建模}：这涉及将离散的用户行为——如屏幕亮度（$B$）、处理器负载（$L$）、网络通信（$S$）、GPS 使用与后台任务——转化为连续或准连续的功率输入信号 $P(t)$~[5,7]。
\item \textbf{物理机理的整合}：模型必须超越简单的能量核算，纳入锂离子电池的固有电化学性质。这一整合聚焦于三个关键维度：温度相关内阻、老化与健康状态（SOH）效应，以及电压崩溃关机机制~[2, 6]。
\item \textbf{预测与定量分析}：模型旨在模拟并近似估计在多样化的初始能量水平与运行场景（如极端寒冷、重度游戏、弱信号通话）下的耗尽时间（TTE），并开展敏感性分析以识别影响电池续航的关键驱动因素~[2,6]。
\item \textbf{验证与实用性}：通过参数估计验证模型以确保其合理性，并最终将研究结果转化为面向智能手机用户与操作系统的可执行策略建议~[8, 13]。
\end{itemize}

\subsection{本文工作}

在本研究中，我们构建了一个完整的连续时间状态空间建模框架，旨在刻画智能手机电池的动态行为，并在多样化的真实使用场景下准确预测耗尽时间（TTE）。我们的方法将多物理场建模与随机仿真及严格的不确定性量化相结合。

我们方法的基础在于构建一个涵盖电化学、热、行为与老化维度的多尺度耦合模型。我们建立了一个二阶 Thevenin 等效电路，并通过具有六个状态变量的连续时间状态空间表示加以形式化，这些状态变量包括荷电状态（SOC）、极化电压以及同步的电池—处理器温度。本阶段的一项关键创新是将 Peukert 效应与温度相关修正整合进 SOC 演化方程，同时引入双节点热反馈模型以捕捉内部产热与处理器负载之间的复杂相互作用。此外，长期退化通过一个结合日历老化与循环老化动力学的混合老化模块加以考虑。

为使这一理论框架可操作，我们实现了一个专门设计的混合数值求解器，用以求解端电压与负载电流之间的非线性代数环。我们的仿真架构采用 5 秒时间步长，以协调迥异的电化学与热时间尺度。模型显式考虑了硬件组件的非线性功耗，例如经 Gamma 校正的屏幕功率与经双曲信号补偿的网络能量。为反映人类行为的不可预测性，我们整合了随机使用模式，采用 Ornstein--Uhlenbeck 过程描述自适应亮度，采用 Bernoulli 过程描述网络尾部能量。该框架通过从日常浏览到极端热环境下高耗电游戏的多场景仿真得到验证。

我们预测的稳健性通过系统性的敏感性与不确定性分析得到进一步强化。我们采用 Morris 筛选法定量参数交互作用，并识别标准与低温条件下不可预测性的主要驱动因素。为评估 TTE 预测相关的极端风险，我们执行 Monte Carlo 仿真（$N=10{,}000$），从而能够计算风险价值（Value-at-Risk, VaR）并评估模型在核心假设压力测试下的表现。这确保我们的建模框架即便在承受显著随机扰动时仍保持稳定。

最后，我们将分析洞见转化为面向终端用户与操作系统的可执行优化策略。这些建议包括面向用户的亮度上限管理与通信批处理，以及面向主动系统级热管理的电压稳定感知动态限流。通过展示模型对无人机（UAVs）与可穿戴技术的可扩展性，我们为更广泛的便携式电子产品谱系中下一代阻抗感知能源管理系统奠定了理论基础。

\section{模型准备}

\subsection{假设与论证}

为在保持计算可处理性的同时准确捕捉锂离子电池在复杂工况下的非线性动力学，我们提出以下核心假设：

\begin{itemize}
\item \textbf{假设 1}：总功耗 $P_{total}(t)$ 视为各独立组件（CPU、屏幕、射频模块等）瞬时功率的算术和，忽略微秒尺度上的电磁耦合与干扰。

\textbf{论证 1}：尽管各组件集成于单一 PCB 上，其电源轨相对独立。在连续时间建模的秒级尺度上，瞬态耦合能量相较运行能量可忽略不计，因此线性叠加是宏观能量流动的稳健近似~[7]。
\item \textbf{假设 2}：电池的电化学动力学由开路电压（OCV）源 $V_{OCV}$、欧姆内阻 $R_0$ 以及两组代表电荷转移与质量扩散的 RC 网络表示。

\textbf{论证 2}：二阶模型在计算复杂度与物理保真度之间取得最优平衡。与一阶模型不同，它能够模拟长时间放电后的电压恢复（松弛效应），同时避免高阶模型带来的过拟合问题~[8, 9]。
\item \textbf{假设 3}：电芯内部温度 $T_{batt}$ 假定为均匀分布，忽略电芯内的径向或空间温度梯度。

\textbf{论证 3}：智能手机电池通常较薄且导热性高，放电倍率很少超过 2C。在这些条件下，毕渥数（Biot number）较低，集总参数（单节点）热模型足以描述电热反馈回路~[2]。
\item \textbf{假设 4}：环境温度 $T_{amb}$ 与网络信号强度 $S(t)$ 等外部变量被视为外生输入，不受设备自身散热或电磁辐射的影响。

\textbf{论证 4}：智能手机耗散的热量相对于环境体积可忽略不计；同样，单台设备的射频功率不会改变全局基站覆盖，从而简化了边界~[6]。
\item \textbf{假设 5}：在单个“耗尽时间”（TTE）预测周期内，健康状态（SOH）假定为常数。老化动力学仅在长期（数月）演化仿真中激活。

\textbf{论证 5}：在单个放电循环（数小时）内的容量损失 $Q_{loss}$ 相对总容量为无穷小量级。将 SOH 视为“慢变量”可显著提升求解器效率，且不损害短期精度。
\end{itemize}

\subsection{符号说明}

本文使用的重要符号列于下。

\begin{description}[leftmargin=8.5em,labelwidth=8em,itemsep=2pt,parsep=0pt,topsep=0.3em]
\item[$\mathrm{SOC}(t)$] 荷电状态（电池电量水平），无量纲。
\item[$\mathrm{SOH}(t)$] 健康状态（$1-Q_{loss}/C_{new}$），无量纲。
\item[$\mathrm{TTE}$] 耗尽时间（剩余可用时间），h/min。
\item[$V_{term}(t)$] 电池端电压，V。
\item[$V_{OCV}(\mathrm{SOC})$] 开路电压（SOC 的函数），V。
\item[$R_0(\mathrm{SOC},T,\mathrm{SOH})$] 欧姆内阻，$\Omega$。
\item[$C_{nom}$] 电池标称容量，Ah。
\item[$I_{batt}(t)$] 电池负载电流，A。
\item[$T_{batt}(t)$] 电池核心温度，$\degC$。
\item[$T_{cpu}(t)$] 处理器温度，$\degC$。
\item[$T_{amb}$] 环境温度，$\degC$。
\item[$B(t)$] 屏幕亮度，无量纲。
\item[$S(t)$] 信号质量因子，无量纲。
\item[$U(t)$] 处理器利用率，无量纲。
\item[$P_{total}(t)$] 总功耗，W。
\item[$\beta(T)$] Peukert 系数（温度相关），无量纲。
\item[$E_{a,R_0}$] 欧姆电阻活化能，J/mol。
\item[$\gamma$] 漏电流温度系数，$1/\degC$。
\end{description}

\section{模型建立}

\subsection{连续时间状态空间表示}

为满足“白箱”机理模型的要求，我们构建了一个多尺度电化学—热—行为—老化耦合动态系统。系统状态向量 $X(t)$ 定义为：
\begin{equation}
X(t)=\left[\mathrm{SOC}(t),\,V_{p1}(t),\,V_{p2}(t),\,T_{batt}(t),\,T_{cpu}(t),\,Q_{loss}(t)\right]^{T}
\end{equation}

各子状态的物理含义定义如下：

\begin{itemize}
\item \textbf{电化学状态}：$\mathrm{SOC}(t)$ 表示荷电状态；$V_{p1}(t)$ 与 $V_{p2}(t)$ 分别表示对应于电荷转移与材料扩散的二阶极化电压~[9]。
\item \textbf{热力学状态}：$T_{batt}(t)$ 与 $T_{cpu}(t)$ 分别表示电池核心温度与处理器温度~[7]。
\item \textbf{老化状态}：$Q_{loss}(t)$ 为慢变变量，表示累积容量损失，用于导出健康状态（SOH）~[3]。
\end{itemize}

为模拟各种用户行为，可控输入向量 $U(t)$ 定义为：
\begin{equation}
U(t)=\left[T_{env}(t),\,\mathrm{RSSI}(t),\,B(t),\,L(t),\,f(t)\right]^{T}
\end{equation}
其中各分量分别表示环境温度、网络信号强度、屏幕亮度、CPU 负载率与时钟频率~[5,6]。

\subsection{电池核心动力学}

\textbf{修正的荷电状态演化}：为捕捉大电流负载下的非线性容量衰减，我们定义纳入 Peukert 效应与热修正的有效放电速率~[8]：
\begin{equation}
\frac{d\mathrm{SOC}(t)}{dt}=-\frac{1}{3600\cdot C_{nom}}\cdot I_{load}(t)\cdot\left(\frac{I_{load}(t)}{I_{ref}}\right)^{\beta(T)-1}
\end{equation}
其中 Peukert 系数 $\beta(T)$ 随温度降低而增大，表征寒冷环境中锂离子迁移受阻所导致的“有效容量崩塌”~[2]。

\textbf{二阶 Thevenin 等效电路}：基于基尔霍夫电压定律（KVL），端电压 $V_{term}$ 由平衡电位、欧姆压降与瞬态极化过电位决定~[2]：
\begin{equation}
V_{term}(t)=V_{OCV}(\mathrm{SOC})-I_{load}(t)R_0(\mathrm{SOC},T,\mathrm{SOH})-V_{p1}(t)-V_{p2}(t)
\end{equation}

\textbf{参数调制与老化耦合}：欧姆内阻 $R_0$ 建模为荷电状态、温度与健康状态的协同函数。老化与热效应之间的耦合由乘性关系表示~[3, 13]：
\begin{equation}
R_0\propto f(\mathrm{SOC})\cdot\exp\!\left(\frac{E_a}{k_BT}\right)\cdot \mathrm{SOH}^{-\gamma}
\end{equation}

\textbf{考虑老化耦合的内阻建模}：欧姆内阻 $R_0$ 由热动力学、电荷分布与结构退化的耦合所表征~[2,3]：
\begin{equation}
R_0=R_{ref}\cdot\exp\!\left[\frac{E_a}{R}\left(\frac{1}{T_{batt}}-\frac{1}{T_{ref}}\right)\right]\cdot\Phi(\mathrm{SOC})\cdot\left(1+\xi_{aging}(1-\mathrm{SOH})^{\lambda}\right)
\end{equation}

这一数学结构表明，随着 SOH 下降，基线电阻增大，并在低温下被 Arrhenius 因子进一步放大。因此，老化设备在寒冷环境中容易出现端电压跌落至截止阈值以下的情况~[2, 3]。

\subsection{负载功率分解与可控输入}

负载电流 $I_{load}(t)$ 是输入向量 $u(t)$ 与系统效率 $\eta$ 的函数：
\begin{equation}
I_{load}(t)=\frac{P_{total}(u(t),T_{cpu})}{\eta\cdot V_{term}(t)}
\end{equation}

总功耗被细分为模块化组件：

\textbf{处理器功耗（$P_{cpu}$）}：计入动态开关与温度相关漏电：$P_{chip}=P_{dyn}(U)+P_{leak0}e^{\gamma(T_{cpu}-T_{ref})}$。

\textbf{显示屏功耗（$P_{screen}$）}：通过 Gamma 非线性建模：$P_{screen}=P_{driver}+k_{br}\cdot B(t)^{2.2}$~[7]。

\textbf{网络功耗（$P_{net}$）}：遵循双曲信号补偿模型：$P_{net}=P_{idle}+\dfrac{P_{tx,base}}{S(t)^{\alpha}}$，其中 $S(t)$ 表示信号强度~[6]。

下图展示了热正反馈循环的示意图：

\subsection{精细双节点热力学反馈}

储能单元与处理单元之间的热交换建模如下~[7]：
\begin{equation}
\begin{cases}
C_{batt}\dfrac{dT_{batt}}{dt}=I_{load}^{2}R_0+\sum\dfrac{V_{pi}^{2}}{R_{pi}}+\dfrac{T_{cpu}-T_{batt}}{R_{int}}-\dfrac{T_{batt}-T_{env}}{R_{ext}}\\[2ex]
C_{cpu}\dfrac{dT_{cpu}}{dt}=P_{chip}-\dfrac{T_{cpu}-T_{batt}}{R_{int}}
\end{cases}
\end{equation}

\subsection{长期老化动力学}

健康状态定义为 $\mathrm{SOH}(t)=1-Q_{loss}(t)/C_{new}$。退化速率由两个非线性分量构成~[3]：

\textbf{日历老化}：SEI 层生长采用基于 Arrhenius 的抛物线律建模，对高 SOC 与热应力敏感：
\begin{equation}
\dot{Q}_{cal}=k_{cal}\cdot\exp\!\left[\frac{-E_{a,cal}}{R}\left(\frac{1}{T_{batt}}-\frac{1}{T_{ref}}\right)\right]\cdot\exp(\alpha\cdot \mathrm{SOC})
\end{equation}

\textbf{循环老化}：机械疲劳与析锂被建模为电流吞吐量与放电深度（DoD）的函数~[3]：
\begin{equation}
\dot{Q}_{cyc}=k_{cyc}\cdot|I_{load}(t)|\cdot(1+\beta\cdot \mathrm{DoD}_{avg})\cdot f(T)^{\gamma}
\end{equation}

其中因子 $f(T)$ 纳入低温运行的惩罚项，以计入析锂导致的加速损伤~[2]。

\subsection{数值实现}

高精度仿真以固定时间步长 $\Delta t=5\,\mathrm{s}$ 执行。这一时间分辨率的选择具有策略性，用以处理电化学极化（秒级）与热传导（分钟级）之间的显著时间尺度差异。该方法确保了迭代非线性代数环的收敛，并维持了系统动力学在多尺度过程中的物理自洽性。

下图展示了典型场景下三条轨迹（SOC、$V_{term}$、$T_{batt}$）的对比图：
\section{耗尽时间预测与场景分析}

\subsection{求解策略与物理负载建模}

锂离子电池系统呈现出一个非线性代数环：端电压决定电流，而电流反过来又影响电压。为求解该问题，我们提出一种混合数值方案~[8]。

\subsubsection{混合求解器算法}

\begin{enumerate}
\item \textbf{电流一致性迭代}：在每个时间步内，对于恒功率负载特性 $P_{total}$，我们执行 5 次快速不动点迭代，以确保欧姆压降、极化电压与负载电流在物理上自洽：
\begin{equation}
V_{term,k+1}^{(t)}=V_{OCV}-\frac{P_{total}^{(t)}}{V_{term,k}^{(t)}}\cdot R_{total}-V_{polarization}^{(t)}
\end{equation}
\item \textbf{操作系统动态干预层}：模型模拟固件级电源管理。若 $T_{batt}<-5\,\degC$ 且端电压 $V_{term}<3.15\,\mathrm{V}$，则触发“低功耗模式”，将功率降低 30\%。
\item \textbf{高精度步进}：维持 $\Delta t=5\,\mathrm{s}$ 的离散步长，以协调二阶极化与分钟级热传导之间的差异。
\end{enumerate}

\subsubsection{关键负载组件的显式非线性建模}

为实现 $P_{total}$ 的精确计算，我们避免简单的线性求和，转而为核心高功耗组件建立非线性物理模型~[7]。

\textbf{1. 屏幕功耗（Gamma 非线性）}

功耗不仅取决于亮度 $B(t)$，还取决于 Gamma 校正系数（通常 $\gamma\approx 2.2$）与平均像素值 $\mu_{pixel}$（特别是对 OLED 面板）：
\begin{equation}
P_{screen}(t)=P_{driver}+k_{br}\cdot B(t)^{\gamma}\cdot\mu_{pixel}(t)
\end{equation}

\textbf{机理}：当亮度从 50\% 调至 100\% 时，功耗并非简单地线性翻倍；相反，它按指数趋势激增约 4.6 倍。

\textbf{2. 网络功耗（信号耦合双曲线）}

我们引入信号质量因子 $S(t)\in(0,1]$ 以模拟射频（RF）放大器的自动增益控制（AGC）~[6]：
\begin{equation}
P_{net}(t)=P_{idle}+\frac{P_{tx,base}}{S(t)^{\alpha}}\cdot I_{transmit}(t)
\end{equation}

\textbf{机理}：随着信号质量 $S(t)$ 退化，发射功率 $I_{transmit}(t)$ 沿双曲线急剧增大。这解释了“信号盲区”（信号死区）中电池续航迅速耗尽这一现象的物理本质。

\subsection{多场景仿真结果与基准验证}

通过标定模型参数，我们获得了在三种典型基准场景下与实测数据高度一致的预测结果~[2,7]：

下图展示了多场景下 TTE 轨迹的叠加图：

注：在“极端寒冷”场景中，失效机理被归类为电压崩溃（电压截止），即低温导致电池内阻显著增大，使得端电压在仍有大量 SOC 剩余时便跌至阈值以下。

\subsection{核心驱动因素分析：电池性能为何“不可预测”}

通过对模型轨迹的深度挖掘，我们识别出导致续航缩短的差异化驱动因素：

\textbf{1. 极端寒冷：内阻分层与“幽灵 SOC”}

在“极端寒冷”场景中，模型揭示了一个反直觉现象：设备在 SOC 仍为 38\% 时便自动关机。

物理驱动因素在于欧姆内阻（$R_0$）与极化电阻（$R_p$）呈现分层 Arrhenius 响应~[9]。在 $-10\,\degC$ 下，$R_0$ 增大约 5--7 倍，造成显著的初始电压跌落。与此同时，$R_p$ 随时间累积，在约 40 分钟内将端电压（$V_{term}$）拉低至 3.0\,V~[2]。因此在低温下，\textbf{端电压}——而非实际能量储备（SOC）——成为 TTE（Time to Empty）的首要瓶颈。

\textbf{2. 重度使用：热电正反馈回路}

在“游戏”场景中，处理器温升导致漏电流随温度指数增长~[7]。

物理驱动因素为 $P_{leakage}\propto e^{\beta T_{cpu}}$~[7]。温度上升使功耗从 4.5\,W 缓慢攀升至 4.8\,W。这种“\textbf{雪球效应}”是高绩效负载下续航加速崩塌的首要原因。

\textbf{3. 屏幕与网络的非线性惩罚}

仿真表明，高屏幕亮度与弱信号导致的续航惩罚远超直觉预期~[6,7]。由于 $P_{screen}\propto B^{\gamma}$，不慎将亮度从 80\% 调至 100\% 会使显示屏额外承受 50\% 的功率惩罚；$P_{network}\propto 1/S^{\alpha}$，弱信号区域射频功耗激增 3--5 倍，迅速耗尽电池。

下图展示了非线性惩罚曲线（Gamma 亮度 + 1/S 信号）：

\subsection{老化效应的影响}

尽管 TTE（Time to Empty）预测主要针对单个放电循环，但根据我们的四层模型架构，SOH（State of Health）的退化对上述预测产生显著的放大效应~[3]：

\begin{itemize}
\item \textbf{内阻耦合机制}：随着循环老化与日历老化推进，内阻 $R_{total}$ 逐渐增大~[3]。对于 SOH $=80\%$ 的“老化”电池，这一变化在室温下可能难以察觉。然而在极端寒冷环境中，基线内阻的增大被低温倍增因子进一步放大。
\item \textbf{预测修正}：模型仿真表明，若将“极端寒冷”场景中的电池替换为老化电池（SOH $=80\%$），关机时的\textbf{剩余 SOC} 将从 38.40\% 恶化至 55\%，TTE 将进一步缩短至 28 分钟。
\end{itemize}

\textbf{结论}：这解释了老旧智能手机在冬季出现“悬崖式”意外关机（电量骤降）背后的物理机制。

下图展示了老化放大效应的对比（SOH=100\% 与 80\%）：

\subsection{模型整体性能评估}

\begin{itemize}
\item \textbf{精度}：对“正常模式”（12.67\,h）的预测与 Carroll 等人（2010）建立的网页浏览基准相比，误差幅度仅为 0.5\%。
\item \textbf{稳健性}：模型成功捕捉了 Peukert 效应~[8] 所引起的低温“容量收缩”与高负载“能量损失”的非线性叠加。这为用户所经历的“SOC 报告不准”（感知到的幽灵电量）提供了定量解释。
\item \textbf{局限性}：当前模型假定环境温度均匀分布。它尚未考虑设备表面局部热点（非均匀热分布）对内阻分布的瞬时影响。
\end{itemize}

\section{敏感性分析}

\subsection{参数交互与非线性探测}

为量化物理常数不确定性对 TTE（Time-to-Empty）的影响，我们采用 Morris 筛选法（$N=512$ 条轨迹）~[12]。我们将“极端寒冷通话”场景（$-10\,\degC$）与“室温”场景（$25\,\degC$）进行比较。通过计算平均基本效应（$\mu^{*}$）与交互强度（$\sigma$），我们识别出模型内部的非线性“杀手”~[12]。

下图展示了 Morris 敏感性雷达图与柱状图：

\textbf{物理洞见：交互作用的本质}

分析表明，$E_a$ 与 $k_{Peukert}$ 之间的交互作用在极端寒冷条件下极为显著（$\sigma$ 值较高）。

\begin{itemize}
\item \textbf{极端寒冷（非线性区制）}：失效模式并非简单的能量耗尽，而是一个恶性正反馈回路：“内阻压降 $\rightarrow$ DC-DC 强制电流增大 $\rightarrow$ Peukert 效应导致额外容量损失。”
\item \textbf{室温（线性区制）}：$E_a$ 的敏感性降至 0.05。在“舒适区”内，物理参数不再是瓶颈，模型回归线性特性。
\end{itemize}

\subsection{复合随机使用中的不确定性量化}

为复现真实续航的“不可预测”本质，我们将三源随机动力学整合进连续时间模型~[5, 6]。

\textbf{随机动力学方程}：

\begin{enumerate}
\item \textbf{亮度随机性（O-U 过程）}：该过程建模用户亮度调节随时间的均值回归行为~[7]。
\begin{equation}
dB(t)=0.1(0.5-B(t))dt+0.05\,dW_t
\end{equation}
（参数基于 Carroll 等人的中位亮度；模拟用户亮度调节的均值回归行为）。
\item \textbf{信号抖动}：基于 Ding 等人（2013）的模型~[6]，信号强度 $S(t)$ 在 $[0.2,0.8]$ 范围内波动，触发非线性发射功率补偿（峰值增益可达 80\%）。
\item \textbf{网络尾部能量}：采用随机 Bernoulli 过程模拟后台任务唤醒~[5]。每次唤醒事件触发 2--4\,s 的高功率状态（450\,mW），模拟 RRC（无线资源控制）状态机的迟滞。
\end{enumerate}

\subsection{TTE 分布与极端风险分析}

通过 $N=10{,}000$ 次 Monte Carlo 仿真，我们量化了 TTE（Time-to-Empty）的不确定性来源：

\textbf{1. 极端寒冷条件（$-10\,\degC$）}

期望 TTE 为 42.4\,min，风险价值（VaR 5\%）为 28.4\,min。物理参数方差的贡献骤增至 38\%（行为随机性占 62\%）。这表明在严酷环境中，电池的“物理体质”（SOH、内阻一致性）成为关键瓶颈。

\textbf{2. 室温条件（$25\,\degC$）对比}

行为随机性贡献了高达 92\% 的方差，而物理参数仅占 8\%。这证实了系统内部存在清晰的主导机制切换。在室温下，用户习惯决定续航；而在极端寒冷中，电池物理属性决定生存。

下图展示了 Monte Carlo TTE 分布的直方图与箱线图：

\subsection{假设压力测试}

我们评估了核心假设被违背时模型预测相对实际结果的绝对偏差，以验证模型的稳健性~[8,9]：

\textbf{关键要点}

\begin{itemize}
\item \textbf{二阶 RC 模型的必要性}：最大偏差（$+19.8\%$）出现在移除扩散极化时，证实浓差极化是“电压受限”关机场景中的主导因素。
\item \textbf{热敏感性}：即便操作系统响应滞后 30 秒，也会造成可测量的续航下降，凸显了软件控制与硬件热极限之间的耦合。
\end{itemize}

\subsection{结论与可扩展性分析}

\textbf{模型稳健性}

尽管输入变量具有高度随机性，模型在整个电池生命周期（SOH $>0.8$）内始终保持优异的物理自洽性。随机过程（如用于亮度的 Ornstein--Uhlenbeck 过程与信号抖动）的整合并未损害底层电化学方程的数值稳定性。

\textbf{设备可扩展性}

该框架的模块化特性使其能够快速适配不同硬件规格。通过将标称容量缩放至 20\,Ah 并重新标定热阻矩阵（$R_{th}$），模型可直接迁移至高耗电动力电池系统，例如无人机（UAVs）所使用者。
\section{系统优化建议}

\subsection{面向智能手机用户的实践}

这些建议的核心在于转变用户行为，使系统远离模型中识别出的高敏感性不稳定区域~[2,6]。

\textbf{1. 亮度的边际效用管理}

\begin{itemize}
\item \textbf{模型洞见}：模型公式 $P_{screen}=P_{driver}+k_{br}\cdot B^{2.2}$ 显示功率与亮度之间呈 Gamma 非线性关系~[7]。
\item \textbf{深度建议}：将手动亮度上限设为 70\%--80\%。
\item \textbf{物理依据}：基于导数 $\frac{dP}{dB}\propto B^{1.2}$，亮度从 80\% 增至 100\% 所带来的功率增量，比从 0\% 增至 20\% 的增量高 5.2 倍。由于人类视觉感知的对数特性，最后 20\% 的亮度在物理上是低效的——即“电池自杀”区。
\end{itemize}

\textbf{2. 弱信号下的“通信批处理”}

\begin{itemize}
\item \textbf{模型洞见}：射频功率 $P_{net}\propto 1/S^{\alpha}$ 与基于 RRC 状态机的“尾部能量”相耦合~[5,6]。
\item \textbf{深度建议}：在信号极弱区域，启用飞行模式或尽量减少高频、小数据包请求（如自动刷新）。
\item \textbf{物理依据}：研究表明，即便是极小的数据包也会触发持续 2--4\,s 的 450\,mW 高功率维持状态。在弱信号下，$\alpha$ 因子的放大效应使频繁的后台心跳将手机困于高功率发射状态，加速 SOC 的崩塌。
\end{itemize}

\textbf{3. 极端寒冷下的大电流规避}

\begin{itemize}
\item \textbf{模型洞见}：内阻 $R_0$ 在极端寒冷下呈 Arrhenius 型指数激增~[2,9]。
\item \textbf{深度建议}：当 $T_{amb}<0\,\degC$ 时，即便 SOC $>20\%$，也应严格避免高瞬态电流应用（如 4K 录像或重度游戏加载）。
\item \textbf{物理依据}：在 $-10\,\degC$ 下，$R_0$ 增大 5--7 倍。大电流 $I$ 会产生巨大的瞬时压降 $\Delta V=I\cdot R_{total}$。即便仍有剩余能量，端电压 $V_{term}$ 也可能瞬间跌至 3.0\,V 物理阈值以下，导致硬关机。
\end{itemize}

基于上述分析，我们在下图中给出推荐行为曲线：

\subsection{操作系统能耗策略}

操作系统应从简单的 SOC 监测转向基于状态向量 $X$ 的主动预测~[13]。

\textbf{1. 基于电压稳定的动态限流（VS-DVFS）}

\begin{itemize}
\item \textbf{改进策略}：引入阻抗感知调度~[13]。当模型检测到 $V_{term}$ 下降，或低 SOC 导致 $R_0$ 漂移时，操作系统应自动限制 CPU 的最大瞬时电流峰值，而非仅仅限制频率。
\item \textbf{物理依据}：传统降频仅考虑“热墙”（温度限制）。然而，我们的二阶 Thevenin 模型表明，极端寒冷与老化导致的电压跌落才是关机的根本原因。通过平滑电流波动，系统可利用电池的极化恢复效应（$V_{p2}$ 的衰减），从而延长有效可用时间。
\end{itemize}

\textbf{2. 泄漏诱发型热失控的主动拦截}

\begin{itemize}
\item \textbf{改进策略}：建立基于 $dT/dt$ 的预测逻辑~[7]。当处理器温升斜率超过预设阈值时，操作系统应在触及热节流墙前 30 秒介入。
\item \textbf{物理依据}：热反馈表现出显著的迟滞性。模型仿真表明，一旦系统进入热电正反馈回路，功耗便呈指数级“滚雪球”。主动拦截切断了正反馈路径，防止系统进入无法通过简单降频逆转的“热失控”状态。
\end{itemize}

\subsection{长期老化与生命周期管理}

\textbf{充电应力控制}：基于老化动力学公式 $\frac{d\mathrm{SOH}}{dt}=f(V,T)$~[3]，建议操作系统引入“峰值电压规避”。在长时间充电期间（如 GPS 导航时），SOC 应锚定在 60\%--80\% 之间，而非维持恒定的 100\%。

\textbf{面向老化的预测修正}：模型仿真表明，老化设备的 TTE（Time-to-Empty）预测变得极不稳定，主要归因于内阻 $R_0$ 的高基线。一旦 SOH $<85\%$，操作系统应放宽 TTE 预测的置信区间，并在 SOC $<30\%$ 时更早发出强制节能提醒。

\subsection{模型推广}

\textbf{1. 无人机}

\begin{itemize}
\item \textbf{推广要点}：无人机运行于 10C--20C 的高倍率放电区~[4]。
\item \textbf{建议}：必须实时监测二阶极化项 $V_{p2}$。由于浓差极化随时间累积，无人机在飞行末段易出现“动力失速”。即便 SOC 看似充足，焦耳热与内阻压降的协同效应也可能触发紧急关机。
\end{itemize}

\textbf{2. 可穿戴设备}

\begin{itemize}
\item \textbf{推广要点}：物理尺寸小导致热阻 $R_{th}$ 极高、热容 $C_{th}$ 极小~[7]。
\item \textbf{建议}：策略应聚焦于\textbf{热质量管理}（Thermal Mass Management）。实施超短脉冲工作模式（占空比 $<5\%$）可使系统利用冷却间隙，防止处理器漏电流超过静态阈值。
\end{itemize}

\subsection{结论}

我们的模型表明，智能手机续航的不确定性并非源于不可预测的用户行为，而是源于电化学系统在非理想环境中的非线性响应。通过从传统的“软件逻辑优化”转向“物理状态感知”，我们能够有效回收原本浪费在内阻压降与热泄漏上的能量——在极端工况下这些损失可达 15\%--25\%。这一范式转变是实现真正长期系统效率与可靠性的关键。

\section{模型的优缺点}

\subsection{模型优点}

\begin{itemize}
\item 模型摒弃了简单的黑箱曲线拟合，采用二阶 Thevenin 电路作为底层逻辑~[8,9]。它成功捕捉了“电压决定电流、电流反作用于电压”的非线性代数环，并能够解释“低温电压击穿”与“热电正反馈”等复杂物理现象~[2,7]。
\item 模型通过不同层次处理从秒级、分钟级到月/年级的跨尺度动力学~[3,9]。这一架构确保了模型在短期 TTE 预测与长期寿命管理两方面的双重有效性。
\item 屏幕功耗采用经 Gamma 校正的非线性模型~[6,7]，网络功耗采用带信号补偿的双曲模型。这种显式的组件建模比传统的线性平均功耗方法更接近真实工况，且实验证明室温下预测误差仅为 0.5\%。
\item 引入 Ornstein--Uhlenbeck（O-U）过程模拟亮度回归，并采用随机 Bernoulli 过程模拟后台尾部能量~[5]。通过 Monte Carlo 仿真，识别出“机制切换”特性：这一深刻洞见表明，室温下续航由用户习惯决定，而极端寒冷条件下生存由物理特性主导~[2,6]。
\item 模型框架不仅适用于手机~[4]。通过调整热阻矩阵与标称容量，它可直接扩展至高倍率放电的无人机与热容极小的可穿戴设备~[7]。
\end{itemize}

\subsection{模型缺点}

\begin{itemize}
\item 我们目前采用双节点集总参数热模型~[7]。在实际设备中，主板上的热点分布极不均匀。忽略空间梯度的简化可能导致对封装极密设备的局部漏电流热失控预测出现偏差。
\item 尽管引入 Arrhenius 方程与抛物线生长律描述老化~[3]，但电化学材料在服役后期的相变或电解液分解具有高度不确定性。在健康状态（SOH）$<75\%$ 后的“性能骤降”区间，模型的预测置信度将下降。
\item 由于非线性代数环的存在，模型需要不动点迭代与高精度步进计算~[8]。若模型在手机端实时运行，其自身的计算开销可能产生不可忽略的二次功耗。
\item 模型将用户行为抽象为随机动力学，但未能准确刻画特定应用在操作系统调度层面的细微差异，这可能在工作条件特定时引入行为预测偏差。
\end{itemize}

\subsection{进一步讨论}

基于当前模型的研究结果，我们认为未来改进与研究的重点应集中于以下方向：

\begin{enumerate}
\item \textbf{引入物理信息神经网络}：为化解计算复杂度与精度之间的权衡，可利用机器学习学习复杂 App 负载的残差项，而将电池核心电化学过程交由物理方程处理~[12]。这种“灰箱”方法既能保持物理可解释性，又能大幅提升对复杂软件行为的预测精度。
\item \textbf{空间精细热建模}：未来模型可从“双节点”结构扩展为“多节点热网络”，甚至纳入二维降阶热图~[7]。这将有助于操作系统在多核调度时避开局部热点，进一步延缓处理器高温引起的漏电流激增。
\item \textbf{云边协同老化修正}：借助手机的网络连接能力，可将数万台设备的充放电数据上传至云端~[6]。通过群体数据分析（Fleet Data Analysis），模型中的经验老化参数可被动态修正~[3]，从而实现面向每位用户、每块电池的“个性化”全生命周期管理。
\item \textbf{对下一代电化学体系的适配}：随着半固态电池或钠离子电池的商业化，本模型的 Thevenin 结构仍然适用~[9]，但其极化项的动态特性将发生显著变化。通过更新参数矩阵，该框架可作为便携式设备能源管理的通用基础。
\end{enumerate}

\subsection{结论}

通过从“软件逻辑优化”转向“物理状态感知”，本模型成功量化了此前被视为“虚标”或“预测不准”的功率崩塌现象。它不仅是预测工具，更为下一代阻抗感知操作系统提供了理论框架与调度策略，因而具有极高的工程实用价值。

\newpage

\section*{参考文献}
\addcontentsline{toc}{section}{参考文献}

\begin{enumerate}[label={[\arabic*]},leftmargin=2.6em,itemsep=0.3em]
\item G. L. Plett, ``Extended Kalman filtering for battery management systems of LiPB-based HEV battery packs Part 1. Background,'' \textit{Journal of Power Sources}, vol. 134, no. 2, pp. 252-261, 2004.
\item K. Smith, C.-Y. Wang, ``Solid-state diffusion limitations on pulse operation of a lithium ion cell for hybrid electric vehicles,'' \textit{Journal of Power Sources}, vol. 161, no. 1, pp. 628-639, 2006.
\item M. Ecker et al., ``Development of a lifetime prediction model for lithium-ion batteries based on extended accelerated aging test data,'' \textit{Journal of Power Sources}, vol. 215, pp. 248-257, 2012.
\item LithiumWerks, ``18650 Lithium Ion Power Cell Nanophosphate\textregistered{} Technology Datasheet (APR18650M1B),'' 2022.
\item J. Huang et al., ``A close examination of performance and power characteristics of 4G LTE networks,'' in \textit{Proc. MobiSys}, Lake Tahoe, 2012.
\item N. Ding et al., ``Characterizing and modeling the impact of wireless signal strength on smartphone battery drain,'' in \textit{Proc. SIGMETRICS}, 2013.
\item A. Carroll, G. Heiser, ``An analysis of power consumption in a smartphone,'' in \textit{Proc. ACM SIGOPS Euromicro}, Nuremberg, Germany, 2010.
\item M. Chen, G. A. Rinc\'on-Mora, ``Accurate electrical battery model capable of predicting runtime and I-V performance,'' \textit{IEEE Transactions on Energy Conversion}, vol. 21, no. 2, pp. 504-511, 2006.
\item M. Doyle, T. Fuller, J. Newman, ``A 1D model of lithium-ion batteries,'' \textit{J. Electrochem. Soc.}, vol. 140, no. 6, pp. 1526-1533, 1993.
\item G. L. Plett, ``Extended Kalman filtering for battery management systems of LiPB-based HEV battery packs. Part 2. Modeling and identification,'' \textit{Journal of Power Sources}, 2004.
\item G. L. Plett, ``Extended Kalman filtering for battery management systems of LiPB-based HEV battery packs. Part 3. State and parameter estimation,'' \textit{Journal of Power Sources}, 2004.
\item A. Pathak et al., ``Fine-grained power modeling for smartphones using system-call tracing,'' in \textit{Proc. EuroSys}, Salzburg, Austria, 2012.
\item O. Bohlen, J. Kowal, D. U. Sauer, ``Impedance-based simulation models for supercapacitors,'' \textit{J. Power Sources}, vol. 172, pp. 468-475, 2007.
\end{enumerate}

\newpage

\section*{AI 使用报告}
\addcontentsline{toc}{section}{AI 使用报告}

依据 2026 MCM/ICM 指南，我们团队使用大语言模型（Large Language Models, LLMs）协助语言优化、数学完善与仿真编码。AI 使用的具体细节分类如下：

\subsection*{1.\ 工具识别与应用}

\subsection*{2.\ 详细使用说明}

\subsubsection*{A.\ 语言翻译与技术写作}

我们论文的初稿，特别是摘要、模型构建与结果分析部分，以中文撰写。我们使用 Gemini 2.5 Pro 执行翻译。为确保学术严谨性，我们应用了特定提示词，如“翻译为专业的 IEEE 期刊风格”。翻译后，我们团队手动修订了 100\% 的技术术语（例如确保“State of Charge”的准确性），以保持上下文一致性。

\subsubsection*{B.\ 电池动力学的数学建模}

为提高电池模型在大电流放电条件下的精度，我们向 Gemini 2.5 Pro 咨询了纳入 Peukert 定律的方法。AI 建议修改荷电状态（SOC）的导数：
\begin{equation}
\frac{dZ(t)}{dt}=-\frac{I(t)^{k}}{Q_n}
\end{equation}
其中 $k$ 表示 Peukert 常数。这一洞见使我们得以改进微分方程组，并在 MATLAB/Python 中开发更稳健的仿真代码。

\subsubsection*{C.\ 随机建模与敏感性分析}

为应对使用模式的不确定性，我们咨询 Grok 4 以修改确定性假设。AI 提出了多层方法：

\begin{itemize}
\item \textbf{随机动力学}：从固定参数转向由布朗运动支配的变量。
\item \textbf{Monte Carlo 仿真}：我们基于 AI 建议的敏感性矩阵结构，实现了 10,000 次迭代的采样框架，以评估模型稳定性。
\end{itemize}

\subsection*{3.\ 人类责任确认}

我们团队确认，核心智力逻辑、问题构思与结果的最终解释完全由人类作者产生。所有 AI 生成的建议，包括公式与代码片段，在纳入最终提交前均经过数学正确性与物理一致性的严格验证。

\end{document}
